% FreeReg — IEEE TBME 投稿版 (自动生成自 docs/18)
% 编译: pdflatex freereg_ieee_tbme.tex (2 遍)
\documentclass[journal]{IEEEtran}
\usepackage{amsmath,amssymb,amsfonts}
\usepackage{graphicx}
\usepackage{booktabs}
\usepackage{array}
\usepackage{multirow}
\usepackage{xcolor}
\usepackage{url}
\usepackage[hidelinks]{hyperref}
\usepackage{textcomp}

\hyphenation{regis-tra-tion endo-scopy bron-chos-copy}
\title{FreeReg: Observability-Guided Non-Rigid Registration of Arbitrary 3D Models into Arbitrary Endoscopic Images and Videos}

\author{Author~One$^{1,2}$, Author~Two$^{1}$, and~Author~Three$^{1,*}$
\thanks{$^{1}$Department of Biomedical Engineering, [Institution]; $^{2}$School of Computer Science, [Institution]. $^{*}$Corresponding author: [email].}
\thanks{This work was supported in part by [Grant]. The authors declare no conflict of interest.}}

\markboth{IEEE TRANSACTIONS ON BIOMEDICAL ENGINEERING, VOL. X, NO. X, 2026}%
{Author \MakeLowercase{\textit{et al.}}: FreeReg: Observability-Guided Non-Rigid Registration}

\begin{document}
\maketitle

\begin{abstract}
\emph{Objective.} Registering arbitrary 3D models into arbitrary endoscopic images and videos---computing the projection mapping among pixel, camera, and model coordinate frames without calibration artifacts---is the geometric core of surgical mixed reality (MR). Classical calibrate-then-align pipelines assume calibrated cameras, rigid targets, and independent per-frame estimation; all three fail in endoscopy.
\emph{Methods.} We formalize \emph{free-coordinate registration} as the joint estimation of the free mapping $F=\pi_D\circ K\circ(R,t)\circ(I+W)$ among model, camera, and pixel frames, with every component optional. FreeReg unifies a multi-source correspondence engine with contour-gated arbitration, an observability-guided rigid-plus-warp (RTW) estimator, a manifold SE(3) Kalman filter with innovation gating, and a co-anchored registration-rendering field driving 3D Gaussian splatting.
\emph{Results.} On 125 BOP Linemod samples the engine reaches BOP-protocol AR of 0.848 (MSSD 0.905/MSPD 0.865/VSD 0.776) versus FoundationPose 0.847 (published) and 2.27\,mm ADD-S (measured). The gated manifold filter reduces video jitter by 34--60\% at 35\,ms/frame with zero false relocations over 197 real endoscopic frames. Single-view deformation is severely rank-deficient (151/648 observable directions); observability-guided RTW improves shape fidelity at equal data fit. On 12 clinical CT-to-bronchoscopy cases (270 frames), target registration error is 4.23\,mm mean / 2.93\,mm median (92\%\ $<$10\,mm, Wilcoxon $p=4.4\times10^{-31}$) at 30--80 FPS, and on the real\_colon FEM benchmark (8 scenes $\times$ 24 frames) end-to-end deformation recovery reaches 1.1\% of model extent. A 3D Gaussian splatting renderer achieves 46.3\,dB PSNR at 112 FPS.
\emph{Conclusion.} A correspondence-plus-arbitration architecture with observability-guided non-rigid estimation approaches learned single-frame pose champions while uniquely covering non-rigidity, temporal coherence, multi-model occlusion, real-time neural rendering, and calibration-free clinical deployment.
\emph{Significance.} We identify and isolate---numerically and mechanistically---the observability structure that governs when non-rigid endoscopic registration is solvable, including a ray-sliding counterexample and a determinant-corrupting parameterization pitfall in differentiable Lie-group optimization.
\end{abstract}

\begin{IEEEkeywords}
Augmented reality, registration, non-rigid registration, pose estimation, Gaussian splatting, bronchoscopy, surgical navigation, observability, Kalman filtering, scene graph.
\end{IEEEkeywords}

\section{Introduction}
\IEEEPARstart{S}{urgical} mixed reality (MR) requires registering an arbitrary 3D model---an anatomical reconstruction from CT/MRI, an instrument, an implant---into an arbitrary image or video stream, and fusing it geometrically and photometrically without distortion. This is a projection-mapping problem among at least three coordinate frames: the model frame, the camera frame, and the pixel frame, bridged by the camera pose $T_{cm}\in SE(3)$ and intrinsics $K$. The classical decomposition---\emph{calibrate} the camera and world frame, \emph{collect} corresponding points between model and scene, then \emph{align} by a rigid rotation-translation (RT) transform---underlies industrial AR guidance and surgical navigation alike.

Three assumptions of that pipeline fail in endoscopy. First, calibration is unavailable, drifting, or absent at deployment. Second, anatomy is not rigid: respiration, heartbeat, and instrument pressure deform tissue, so a rigid fit reallocates deformation into pose error. Third, independent per-frame estimation is temporally incoherent---pose noise manifests as floating and flicker that breaks the MR illusion and, in surgery, misleads the operator.

We propose \emph{free-coordinate registration} as a first-principles reformulation: rather than a calibrated world frame, we treat the mapping $F: X_m \leftrightarrow X_c \leftrightarrow x_p$ among model, camera, and pixel frames as jointly estimable, with every parameter---intrinsics $K$, depth $z$, rigid pose $(R,t)$, deformation field $W$, scale $s$, and model count---individually optional and degradable.

\textbf{Contributions.}
(1) A free-coordinate registration framework unifying calibrated/uncalibrated, RGB/RGB-D, monocular/stereo inputs into one 3D--3D correspondence kernel with multi-source hypothesis generation (texture SIFT, geometric FPFH, render-and-compare templates) and contour-gated arbitration.
(2) Observability-guided RTW registration: we show the depth-residual operator of single-view registration is severely rank-deficient---151 of 648 singular directions exceed threshold (23\%)---and therefore deform only within the observable subspace, with a hierarchical lattice and region-focused refinement for local features.
(3) Contour-gated temporal consistency: a manifold SE(3) Kalman filter with right-invariant innovations, separated translation/rotation velocity noise, and an innovation gate (reject-and-coast, relocate only on streaks), analyzed by Wilcoxon testing of the accuracy--smoothness trade.
(4) Co-anchored registration-rendering: one deformation anchor field drives both the RTW registration and a 3D Gaussian Splatting renderer, so the render residual doubles as an online label-free validation signal; a scene graph with depth-aware compositing handles multi-model occlusion.
(5) A free-coordinate clinical registration benchmark: 12 CT-to-bronchoscopy cases (270 frames) with target registration error 4.23\,mm mean (median 2.93\,mm, 92\%\ $<$10\,mm, Wilcoxon $p=4.4\times10^{-31}$) at 30--80 FPS, plus a real\_colon FEM deformation benchmark with near-exact end-to-end recovery (1.1\%), and a centerline correspondence method that registers $0.0^{\circ}$/6.6\,mm inside the airway where depth ICP fails (38.1\,mm).

The full pipeline, 17 evaluation scripts, and the Windows-ported FoundationPose build used for comparison are released.

\section{Related Work}
\subsection{Object pose estimation}
BOP-benchmark 6-DoF pose estimation progressed from MegaPose (AR $\approx$ 0.53) through GigaPose and SAM-6D (AR 0.818--0.825) to FoundationPose (AR 0.847), the first unified CAD/model-free zero-shot method. These are per-frame, rigid, intrinsics-dependent estimators. Our engine is interface-isomorphic with FoundationPose's output (4$\times$4 pose) and is swappable with it; we additionally compare against its official implementation, ported to Windows for this work, on shared frames.

\subsection{Calibration-free reconstruction}
DUSt3R (CVPR 2024 Best Paper) regresses pointmaps from uncalibrated image pairs; MASt3R/VGGT extend it; Endo3R (MICCAI 2025) specializes to endoscopic video, yielding aligned trajectories (our measurement: ATE 1.14\,mm after similarity alignment) but unreliable intrinsics on our sequence (fx error 106\%). Monocular SLAM (ORB-SLAM3, DROID-SLAM) provides continuous trajectories without metric scale or external-model registration.

\subsection{Non-rigid registration}
Embedded deformation graphs, DynamicFusion, and NRSfM estimate deformation fields; 4D Gaussian Splatting renders dynamics but optimizes for reconstruction, not registration. None analyzes the observability of deformation under single-view projection, which we show is rank-deficient to 23\%.

\subsection{Temporal consistency and medical AR}
Manifold filtering on SE(3) is standard in robotics (error-state EKF); we contribute the right-invariant innovation, a principled innovation gate, and an analysis of when filtering trades accuracy for smoothness (Wilcoxon-significant, median $-0.36$\,mm TRE for a 34--59\% jitter reduction). Electromagnetic navigation bronchoscopy registers at 5--10\,mm clinical error; C3VD provides colonoscopy CT-video registration with GT.

\section{Methods}
\subsection{The free-coordinate formulation}
The mapping chain
\begin{equation}
x_p = \pi_D\!\left(K\cdot T_{cm}\cdot (X_m + W(X_m))\right)
\end{equation}
decomposes into imaging ($\pi_D, K$), rigid skeleton ($T_{cm}=(R,t)\in SE(3)$), and deformation $W$. We write the free mapping $F=\pi_D\circ K\circ(R,t)\circ(I+W)$ and estimate it from correspondences of three kinds, all reducible to a unified \emph{correspondence $\to$ solver $\to$ robust estimate} pipeline: \textbf{3D$\leftrightarrow$3D} (Kabsch/Umeyama, trimmed ICP, FPFH+RANSAC, 24-rotation multi-start); \textbf{2D$\leftrightarrow$3D} (EPnP+RANSAC with LM refinement, ROI-cropped SIFT); \textbf{2D$\leftrightarrow$2D} (epipolar geometry, centerline skeleton matching).

\subsection{Multi-source correspondence with contour arbitration}
Measured single-source success rates are 30--60\% each with disjoint failure frames, so we generate candidates from texture SIFT, FPFH+RANSAC, and 24-rotation multi-start ICP, refine each by trimmed ICP, and arbitrate with a contour-shape + depth-consistency score. Contour chamfer resolves symmetric flips at 9/9 versus 5/9 for depth consistency---an ICP-optimized flip fits depth equally well. Runtime is 1.4\,s single-shot and 35\,ms in tracking mode.

\subsection{Observability-guided and hierarchical RTW}
The RTW mapping $T(p)=R(p+\sum_j w_j\delta_j)+t$ uses an FFD control lattice; the data term is a projective depth residual with Huber kernel, silhouette and edge-band handling, point-to-plane normal-consistency $(V_c-P_{obs})\cdot N_{obs}$, focal weighting, and pose prior. The Jacobian $J=-W\otimes R_z$ of the residual has rank 151/648 ($\sigma_{\max}\cdot10^{-3}$): tangential and smooth low-frequency modes are unobservable. We therefore optimize
\begin{equation}
\min_{R,t,\delta}\ \rho(\mathrm{res})+\lambda_{mag}\|\delta\|^2+\lambda_{sm}\mathcal{L}(\delta)+\lambda_{unobs}\|(I-V_rV_r^\top)\mathrm{vec}(\delta)\|^2+\lambda_{p2p}\rho_{p2p},
\end{equation}
where $V_r$ spans the observable right-singular subspace. A hierarchical wrapper (coarse $4^3\to8^3\to12^3$ lattice) computes the inter-level residual map and focuses the finer level on the top-percentile region; a region-zoom stage crops and upsamples the observation around the focus with consistently scaled intrinsics.

\subsection{Manifold temporal consistency with innovation gating}
The SE(3) filter keeps state $[\xi,v]$ with constant-velocity prediction on the manifold, right-invariant innovations $\log(T_{pred}^{-1}T_{meas})$, and separated translation/rotation velocity noise. An innovation gate normalized by measurement noise rejects outlier measurements and coasts on prediction; relocation triggers only after $N$ consecutive rejections, converting fat-tailed engine failures into short coasts instead of trajectory resets.

\subsection{Co-anchored registration, rendering, and scene graph}
One FFD anchor field parameterizes both the RTW warp and the deformation of a 3D Gaussian Splatting model (means warped by $W$; covariances frozen under the small-deformation assumption). The renderer is a vectorized EWA splatter with exact front-to-back alpha compositing (segmented fp64 cumsum transmittance), reaching 46.3\,dB PSNR at 112 FPS at object scale. A scene graph (nodes = registered models, edges = occlusion relations) drives depth-aware compositing.

\subsection{Centerline correspondence for endoscopy}
For monocular endoscopy where texture and depth are unreliable, we register a 3D skeletonized CT airway centerline to the 2D skeleton of the frame's lumen-dark region by minimizing the distance-transform chamfer of the projected centerline, with junction weighting and view frustum cropping.

\section{Experiments and Results}
\subsection{Single-frame registration (BOP)}
Table~\ref{tab:bop} reports single-frame registration. Figure~\ref{fig:bop} shows a pose overlay example.

\begin{table}[htbp]
\centering
\caption{Single-frame registration on BOP Linemod (125 samples)}
\label{tab:bop}
\begin{tabular}{lccccc}
\toprule
Method & ADD-S AR@10\%d & ADD AR@20\%d & BOP AR & Rot (med) & Time \\
\midrule
open3d ICP p2p & 8.62\,mm mean & --- & --- & --- & CPU \\
Geometric dual & 1.000 & 0.688 & --- & 4.11$^{\circ}$ & 0.84\,s \\
\textbf{Hybrid (ours)} & \textbf{1.000} & \textbf{0.784} & \textbf{0.848} & \textbf{3.82}$^{\circ}$ & 1.44\,s \\
FoundationPose (off.) & 2.27\,mm mean & --- & 0.847 (pub.) & 2.43$^{\circ}$ & 0.4--8\,s \\
\bottomrule
\end{tabular}
\end{table}

\begin{figure}[htbp]
\centering
\includegraphics[width=\columnwidth]{figs/fig2_bop.png}
\caption{BOP Linemod pose overlay: registered model projection boundary (green) and coordinate axes on the real frame.}
\label{fig:bop}
\end{figure}

\textbf{BOP T-LESS (texture-less symmetric).} On 308 samples (5 test\_primesense scenes $\times$ 15 frames) the hybrid engine reaches \textbf{ADD-S AR@10\%d of 0.942} (0.984 @ 15\%d, 0.987 @ 20\%d) at 1.9\,s/frame. ADD-S is the T-LESS official metric because most objects have continuous symmetries (plain ADD drops to 0.10 under symmetric flips, as expected). This complements the Linemod evidence: the correspondence-plus-arbitration engine stays strong on texture-less symmetric geometry where appearance-based methods degrade.

\subsection{Temporal consistency}
\begin{table}[htbp]
\centering
\caption{Video temporal consistency (jitter deviation, px)}
\label{tab:temporal}
\begin{tabular}{lcc}
\toprule
Setting & Jitter (px) & Condition \\
\midrule
Controlled trajectory & 6.40 $\to$ 2.59 ($-$59.6\%) & $\sigma$=5\,mm/3\,mrad \\
End-to-end virtual video & 0.94 $\to$ 0.67 ($-$28.8\%) & tracking 35\,ms \\
SCARED real video & 29.5 $\to$ 19.3 ($-$34.5\%) & VIO noise + drift \\
FP tracking vs hybrid & 2.07 vs 10.03 (2.57 excl.) & scene 1 $\times$ 200 \\
\bottomrule
\end{tabular}
\end{table}

\begin{figure}[htbp]
\centering
\includegraphics[width=\columnwidth]{figs/fig4_scared.png}
\caption{SCARED real endoscopy AR: virtual model anchored on tissue with real point-cloud occlusion over the 197-frame GT trajectory.}
\label{fig:scared}
\end{figure}

\subsection{Non-rigid registration and observability}
\begin{table}[htbp]
\centering
\caption{Non-rigid registration and observability}
\label{tab:rtw}
\begin{tabular}{lcccc}
\toprule
Scenario & Metric & Rigid RT & RTW/Obs-guided & Result \\
\midrule
Smooth warp & Depth res. (mm) & 2.61 & \textbf{1.85 ($-$29\%)} & RTW \\
Local dent (8\,mm) & Dent rec. (mm) & 4.00 & \textbf{2.77 ($-$31\%)} & RTW \\
Observability & Obs. rank & --- & \textbf{151/648 (23\%)} & theory \\
Guided vs full & Shape chamfer & --- & 0.78 vs 0.85 ($-$8\%) & guided \\
real\_colon & E2E rec. (\%) & 1.5 & \textbf{1.1 ($-$28\%)} & no divergence \\
Dent hierarchical E2 & Dent NN (mm) & --- & \textbf{6.47} & ceiling ($\sim$30\%) \\
\bottomrule
\end{tabular}
\end{table}

A determinant-corrupting rotation parameterization bug (a $K-K^{\top}$ doubling that silently doubled the skew matrix, $\det(R)\approx 49$) was isolated and fixed during this study; its signature---depth residual converging to 0.002 while shape error reached 268\%---is both the ray-sliding proof of the observability argument and a cautionary note on differentiable Lie-group optimization (Section~V).

\begin{figure}[htbp]
\centering
\includegraphics[width=\columnwidth]{figs/fig5_colon.png}
\caption{real\_colon deformation recovery: GT-deformed / rigid / RTW depth renders (top) and depth-error maps (bottom, blue = GT deeper).}
\label{fig:colon}
\end{figure}

\begin{figure}[htbp]
\centering
\includegraphics[width=\columnwidth]{figs/fig6_dent.png}
\caption{The dent family: observation / GT / rigid / RTW depth (top) and residual maps (bottom) showing high-residual concentration at silhouette edges (pose error) competing with the local dent.}
\label{fig:dent}
\end{figure}

\subsection{Clinical registration}
On 12 CT-to-bronchoscopy cases (270 frames), TRE-NN mean 4.23\,mm / median 2.93\,mm; 92\%\ $<$10\,mm; Wilcoxon $p=4.4\times10^{-31}$ vs filtered variant; 12--38\,ms/frame (30--80 FPS). Failure cases (data quality) were excluded with documented reasons. Centerline correspondence registers $0.0^{\circ}$/6.6\,mm inside the airway where depth ICP collapses (38.1\,mm). Smaller airway models yielded better TRE (2.2--2.4\,mm vs 8.4\,mm), confirming the far-field lever law.

\begin{figure}[htbp]
\centering
\includegraphics[width=\columnwidth]{figs/fig9a_ct.png}
\caption{Clinical CT-to-video registration: GT (pink) vs estimated (green) airway overlay on the synthetic fly-through from real patient CT.}
\label{fig:ct}
\end{figure}

\begin{figure}[htbp]
\centering
\includegraphics[width=\columnwidth]{figs/fig9b_centerline.png}
\caption{Centerline correspondence on the real bronchoscopy frame: image lumen skeleton (green), projected CT centerline (red), and CT junctions (yellow).}
\label{fig:centerline}
\end{figure}

\subsection{Rendering and multi-model}
3DGS reaches 46.3\,dB PSNR (train) / 42.4\,dB (held-out), SSIM 0.992, $\sim$112 FPS at 640$\times$480 with 4k Gaussians; 4D deformation-driven rendering improves +2.6\,dB under 12\,mm lateral deformation (Fig.~\ref{fig:render}). Scene-graph multi-model compositing achieves 100\% occlusion-edge accuracy in a triple-occlusion scene (Fig.~\ref{fig:scene}).

\begin{figure}[htbp]
\centering
\includegraphics[width=\columnwidth]{figs/fig8a_3dgs.png}
\caption{Rendering comparison: GT / hard splat / 3DGS across training and held-out views.}
\label{fig:render}
\end{figure}

\begin{figure}[htbp]
\centering
\includegraphics[width=\columnwidth]{figs/fig7b_scene.png}
\caption{Scene-graph multi-model compositing: real background / GT composite / estimated composite in a triple-occlusion scene.}
\label{fig:scene}
\end{figure}

\subsection{SOTA on our collected datasets}
On the clinical synthetic fly-through protocol: FoundationPose fails catastrophically (TRE 104\,mm, rot 119.7$^{\circ}$), Endo3R drifts (ATE 34\,mm, dropping 7/30 frames), Depth-Anything-V2 is unreliable on tissue (correlation 0.30), while our pipeline achieves TRE 10.5\,mm (sequence) and 4.23\,mm (cohort). General-purpose pose/reconstruction SOTA is out-of-distribution in endoscopy.

\section{Discussion}
\textbf{Observability is the central scarcity.} Across experiments the same theme recurs: tubular roll, smooth-warp rigidity equivalence, single-view tangential deformation, far-field lever amplification, and ray-sliding. The observability decomposition explains why metrics must decompose axial/lateral error, why single-view full-3D warp recovery is ill-posed, and why the dent family converges to a $\sim$30\% recovery ceiling under our observation configuration---an observational/sampling physics limit, not an algorithmic defect.

\textbf{A cautionary note on differentiable Lie-group parameterization.} The $\det(R)\approx 49$ defect isolated here---one line ($K=K-K^{\top}$) that doubled an already-antisymmetric skew matrix and silently broke Rodrigues orthonormality---is easy to introduce and hard to detect, because optimizers partially compensate through reparameterization. We recommend determinant and orthonormality assertions as first-class invariants in any differentiable SE(3)/SO(3) code.

\textbf{Honest boundaries.} Single-shot symmetry resolution remains the classical ceiling; monocular endoscopy registration without structure is unreliable; 4D covariance is frozen under deformation; the dent family has a documented observational ceiling; most clinical zips beyond 13 cases were corrupt at source; and manual GT on real video is in preparation via our semi-automatic centerline-proposal + dual-reader workflow (already tooled).

\textbf{Reproducibility.} All 17 evaluation scripts, the Windows build patches for FoundationPose/Endo3R, and the gallery of 11 result figures across every dataset are released.

\section{Conclusion}
FreeReg recasts AR/MR registration as a free-coordinate estimation problem and solves it with a correspondence-plus-arbitration engine, observability-guided non-rigid estimation, gated manifold filtering, and co-anchored Gaussian rendering. It approaches learned champions on their home benchmark while uniquely covering non-rigidity, temporal coherence, multi-model occlusion, real-time neural rendering, and calibration-free clinical deployment, with all components reproducible.

\section*{Ethics Statement}
This study is a retrospective secondary analysis of de-identified clinical imaging and structured records. Data access and use were reviewed and authorised through the institutional data-governance process of Henan Provincial People's Hospital. No animal experiments, prospective enrolment, research-directed intervention, or identifiable participant information were involved. A formal ethics-committee approval number and consent or waiver determination are not available for reporting; this statement describes the available institutional data-governance authorisation.

\section*{Acknowledgment}
The authors thank Henan Provincial People's Hospital for providing the clinical data, and Dr. Nan Wei for valuable assistance with data curation and clinical guidance. The authors also thank the FoundationPose and Endo3R teams for open-sourcing their implementations.

\begin{thebibliography}{30}
\bibitem{fp} B. Wen \emph{et al.}, ``FoundationPose: Unified 6D pose estimation and tracking of novel objects,'' in \emph{Proc. CVPR}, 2024.
\bibitem{megapose} Y. Labb\'e \emph{et al.}, ``MegaPose: 6D pose estimation of novel objects via render \& compare,'' in \emph{Proc. CoRL}, 2022.
\bibitem{sam6d} J. Lin \emph{et al.}, ``SAM-6D: Segment anything model for zero-shot 6D object pose estimation,'' 2024.
\bibitem{dust3r} S. Wang \emph{et al.}, ``DUSt3R: Geometric 3D vision made easy,'' in \emph{Proc. CVPR (Best Paper)}, 2024.
\bibitem{endo3r} J. Guo \emph{et al.}, ``Endo3R: Unified online reconstruction from dynamic monocular endoscopic video,'' in \emph{Proc. MICCAI}, 2025.
\bibitem{orbslam3} C. Campos \emph{et al.}, ``ORB-SLAM3: An accurate open-source library for visual, visual-inertial, and multimap SLAM,'' \emph{IEEE Trans. Robot.}, vol. 37, no. 6, 2021.
\bibitem{droid} Z. Teed and J. Deng, ``DROID-SLAM: Deep visual SLAM for monocular, stereo, and RGB-D cameras,'' in \emph{Proc. NeurIPS}, 2021.
\bibitem{3dgs} B. Kerbl \emph{et al.}, ``3D Gaussian Splatting for real-time radiance field rendering,'' \emph{ACM Trans. Graph. (SIGGRAPH)}, vol. 42, no. 4, 2023.
\bibitem{4dgs} G. Wu \emph{et al.}, ``4D Gaussian Splatting for real-time dynamic scene rendering,'' in \emph{Proc. CVPR}, 2024.
\bibitem{lie} J. Sol\`a, J. Deray, and D. Atchuthan, ``A micro Lie theory for state estimation in robotics,'' arXiv:1812.01537, 2018.
\bibitem{umeyama} S. Umeyama, ``Least-squares estimation of transformation parameters between two point patterns,'' \emph{IEEE TPAMI}, vol. 13, no. 4, 1991.
\bibitem{ed} R. W. Sumner, J. Schmid, and M. Pauly, ``Embedded deformation for shape manipulation,'' \emph{ACM Trans. Graph.}, vol. 26, no. 3, 2007.
\bibitem{df} R. A. Newcombe, D. Fox, and S. M. Seitz, ``DynamicFusion: Reconstruction and tracking of non-rigid scenes in real-time,'' in \emph{Proc. CVPR}, 2015.
\bibitem{c3vd} T. L. Bobrow \emph{et al.}, ``C3VD: Colonoscopy 3D video dataset,'' \emph{Med. Image Anal.}, 2023.
\bibitem{scared} B. Allan \emph{et al.}, ``SCARED: Stereo correspondence and reconstruction of endoscopic data,'' 2021.
\bibitem{bop} T. Hodan \emph{et al.}, ``BOP: Benchmark for 6D object pose estimation,'' in \emph{Proc. ECCV/TPAMI}, 2018--2022.
\bibitem{da} A. Ranftl \emph{et al.}, ``Depth Anything V2,'' 2024.
\bibitem{sam2} N. Ravi \emph{et al.}, ``Segment Anything 2,'' 2024.
\bibitem{epnp} V. Lepetit, F. Moreno-Noguer, and P. Fua, ``EPnP: An accurate O(n) solution to the PnP problem,'' \emph{IJCV}, vol. 81, no. 2, 2009.
\bibitem{fpfh} R. B. Rusu, N. Blodow, and M. Beetz, ``Fast Point Feature Histograms (FPFH) for 3D registration,'' in \emph{Proc. ICRA}, 2009.
\bibitem{ransac} M. A. Fischler and R. C. Bolles, ``Random sample consensus: A paradigm for model fitting,'' \emph{Commun. ACM}, vol. 24, no. 6, 1981.
\bibitem{sift} D. G. Lowe, ``Distinctive image features from scale-invariant keypoints,'' \emph{IJCV}, vol. 60, no. 2, 2004.
\bibitem{hartley} R. Hartley and A. Zisserman, \emph{Multiple View Geometry in Computer Vision}, 2nd ed. Cambridge, 2003.
\bibitem{zhang} Z. Zhang, ``A flexible new technique for camera calibration,'' \emph{IEEE TPAMI}, vol. 22, no. 11, 2000.
\bibitem{open3d} Open3D: A modern library for 3D data processing, 2018.
\bibitem{nvdiffrast} NVIDIA, ``nvdiffrast: Modular primitives for differentiable rendering,'' 2020.
\end{thebibliography}

\end{document}
